\documentclass{subfiles}
\begin{document}\label{sec:rb-trees}
    A special variant of \gls{bst}, RB trees allow to implement the operations of 
        \lstinline{Search(S, x)}, \lstinline{Insert(T, x)} and \lstinline{Delete(T, x)} in 
        \(\ofO{\log n}\) time. 

    Formally speaking, RB trees are binary search trees satisfying the following properties:
    \begin{enumerate}
        \item \label[Property]{prop:rb-prop-1} Each internal node is either red or black.
        \item \label[Property]{prop:rb-prop-2} The root is black.
        \item \label[Property]{prop:rb-prop-3} Each leaf is black.
        \item \label[Property]{prop:rb-prop-4} Each red node has only black children.
        \item \label[Property]{prop:rb-prop-5} Any simple path from a node to a descendant leaf,
            contains the same number of black nodes.
    \end{enumerate}

    \Cref{prop:rb-prop-5} implicitly defines a measure to which, for a given node \(x\),
        we refer to as the black height \(\bh\). That is,
        \(\bh\) is the number of black nodes in the path from \(x\) to any of its descendant leaf.
        The following result easily follows.
        \begin{lemma}
            If \(T\) is a Red-Black tree with \(n\) nodes, 
                then \(T\) has height at most \(2 \log(n + 1)\).
        \end{lemma}
        \begin{proof}
            Follows immediately by a simple induction on \(n\).
        \end{proof}

    \paragraph{Rotations}
    \subfile{../paragraphs/RB rotations}

    \paragraph{Insertion}
    \subfile{../paragraphs/RB insertion}

    \paragraph{Deletion}
    \subfile{../paragraphs/RB deletion}
    \clearpage
\end{document}
